\documentclass[11pt]{article}
\usepackage[margin=0.8in]{geometry}
\usepackage{booktabs,tabularx,xurl,hyperref,amsmath}
\hypersetup{hidelinks}\setlength{\parindent}{0pt}\setlength{\parskip}{6pt}
\title{Coarsened Sparse SQP--RTI Deployment Validation}
\author{Collision Avoidance Research}\date{November 4, 2026}
\begin{document}\maketitle
\section*{Implemented result}
The active raw-sensor controller now solves one sparse multiple-shooting QP
per call. It optimizes shared actuator and tanh-memory coefficients, then commits
the first policy directly. Whole-horizon coefficient finite differences,
onlinear backtracking, nonlinear candidate acceptance, and iterative cold-start
feasibility restoration were removed from the deployment path. The exact
robust theorem remains in \path{controller/CONTROLLER_THEORY.md}.

All 225 focused MATLAB tests passed. Five final closed-loop cases each completed
200 holds. Two of the five initializations exceeded 50 ms; all subsequent
995 continuation calls stayed below 50 ms, with a maximum of 43.629 ms on
this workstation. Initialization must finish before the timed control loop. This is measured performance, not a worst-case
execution-time or universal safety guarantee. Both original 25 m/s challenge
cases remain outside the successful QP-admission domain and issued no command.

\section*{Timing and physical replay}
Hardware: AMD Ryzen 7 7800X3D (8 cores, 16 threads); Linux 7.3.1-1-cachyos;
MATLAB R2026a Update 3 with Optimization Toolbox. Four discarded virtual startup calls
warm initialization and continuation in MATLAB and the solver before each main/dropout run. The timing campaign runs in a separate fresh MATLAB process after the focused test process exits. Timing includes the
entire controller API: information update, nominal prediction, local derivatives,
QP and command bookkeeping. Independent tight-tolerance \texttt{ode45} plant
replay, sensor synthesis, and JSON export are outside that controller timing.
They are simulation infrastructure rather than controller work.

\begin{center}\small
\begin{tabular}{@{}rlrrrrrr@{}}\toprule
Speed & Case & Holds & Median & P95 & Max & Misses & Min gap\\
(m/s) & & & (ms) & (ms) & (ms) & $>50$ ms & (m)\\\midrule
1.5 & braking & 200 & 24.81 & 34.73 & 36.66 & 0 & 25.200 \\
1.5 & stopped & 200 & 24.54 & 32.64 & 102.26 & 1 & 28.505 \\
8 & braking & 200 & 11.52 & 38.40 & 41.34 & 0 & 0.300 \\
8 & stopped & 200 & 17.12 & 40.66 & 53.06 & 1 & 0.300 \\
8 & dropout & 200 & 11.89 & 38.47 & 42.41 & 0 & 0.300 \\
\bottomrule\end{tabular}\end{center}

A fresh MATLAB process with explicit warmup disabled recorded controller calls
of 903.8, 87.3, 65.0 ms. Startup must precede the timed loop. Earlier exploratory
runs also showed occasional desktop latency outliers above 50 ms; the
final continuation result does not rule out future outliers. No real-time operating
system, deadline interrupt or actuator latency qualification is claimed.

\section*{Default deployment and validation scope}
The fixed horizon has four 50 ms stages, one memory coordinate and two future
sensor substeps per hold. Each retained physical hypothesis has one zero-noise
future case; the experiments start with three correlated hypotheses. Actual
12.5 ms sensor readings all condition the posterior. The finite prior contains
the actual initial member by construction; it is not a sensor-derived enclosure.
Scenario seed is 20261104. Dropout holds are 2, 3, 20 and 21. Each hold has 41
independent replay clearance checks, with \texttt{ode45} relative tolerance
$10^{-10}$ and absolute tolerance $10^{-12}$.

\section*{Computational and algorithmic changes}
Each stage perturbs only six physical coordinates and two actuator coordinates.
Analytic derivatives handle memory, shared coefficients, previous input and
geometric witnesses. Shared policies couple sparse scenario state chains;
root memory is fixed and all root hypotheses receive the same action. The
shooting equalities, inequalities and regularized Gauss--Newton Hessian remain
sparse through \texttt{quadprog}. Every successful recorded call uses one QP.
The linearized shooting forecast is labeled explicitly; auxiliary physical
nodes and sensor blocks are nominal, not verified trial trajectories.

Collision planes use start/mid/end samples and a 0.2 m engineering reserve in
addition to the 0.1 m requested clearance. The approximate terminal plane adds
constant-deceleration travel and a fading-brake travel term. It does not enforce
membership in the theoretical circular invariant set. Command magnitude and
slew rows are hard QP inequalities; state and geometry rows allow half the
configured $10^{-5}$ numerical tolerance. No nonlinear verification transfers
these QP properties to the physical system.

Earlier diagnostic terminal variants failed during an 8 m/s stop: the circular
reserve left a roughly $10^{-6}$ local feasibility discrepancy at rest, and a
constant-deceleration plane reserve omitted fading-brake travel. Adding the
fading term completed the 200-hold stop in the final tests. These diagnostics
were not accepted as successful final validation.

\section*{Checks, limitations and reproducibility}
The eight focused classes cover joint causality, complete raw-block conditioning,
parameter persistence, actuator timing, one-QP execution, numerical shooting
tangents against independent policy perturbations, configuration, source budget,
tires, road load, physical terminal calculations and input geometry. Full
repository tests were not run. Code Analyzer outputs are preserved; the missing
configured R2025b settings file caused MATLAB to use defaults. Analyzer findings
include style/performance suggestions and are not reported as warning-free.

The 25 m/s stopped/braking-lead cases failed their single local QP with exit
flag $-2$. This is not proof of physical unavoidability. The short horizon,
restricted backup tail and approximate terminal reserve limit feasibility.
No unsafe-initialization recovery, arbitrary continuing/restarting target,
multiple-target, hard-road, or universal robust-safety claim is made. A QP error
issues no new command but does not itself provide a safe continuation policy.

Exact commands, all raw traces and per-call timing components, focused test
results and analyzer diagnostics are under \path{report/SPARSE_RTI_20261104/}.
The preceding 9--10 second measurements used different dimensions and a different
algorithm; they are not a controlled speedup experiment. The revised measured
latency and its full timing distribution are the deployment result.
\end{document}
